\section{Transverse density and isolated-pair interactions}
\label{sec:gt-transverse}

The complete horizontal count in Section~\ref{sec:gt-height} also bounds
the number of terms in broad transverse observables. Their evolution
retains a dissipative contribution from conjugate zeros. This gives
a heat-time estimate in addition to the bounds available from a single
slice. Write all zero sums with multiplicity, use physical time
\(0\le t\le1/2\), and use the canonical strip \(|\Im\rho|\le1\)
throughout this interval.
The growth proof preceding Proposition~\ref{gt-exterior-proposition}
applies on this same slab with \(u^2/2\) in the source envelope.

\subsection{A convex balance with conjugate dissipation}

Write \(\mathcal L_t(z)=H_t'(z)/H_t(z)\), and take \(C_{\rm u}\ge1\)
as a common constant in \eqref{gt-height-unit-count} and
\begin{equation}
 |\mathcal L_t(u+4\ii)|\le C_{\rm u}\log(2+|u|)
 \qquad(u\in\R,\ 0\le t\le1/2).
 \label{gt-safe-line}
\end{equation}
The latter estimate for positive time is the upper-line logarithmic
derivative bound in the proof of Theorem~1.5(iv) of Polymath
\cite[\S9, following (105), using Proposition~9.1(ii)]{polymath}.
At time zero, the source formula gives locally uniform convergence of
\(H_t,H_t'\). The line is zero-free throughout the closed heat slab,
so for each \(u\) the same constant follows on taking \(t\downarrow0\).
Reflection and compactness cover negative and bounded \(u\).

Fix \(\chi\in C_c^2(\R)\), supported in \([-1,1]\), with
\(0\le\chi\le1\) and \(\chi=1\) on \([-1/2,1/2]\). For \(U\ge8\) set
\[
 \psi_U(x)=\chi((x-3U/2)/U),\qquad
 \mathcal C_U=\{j:U/4\le x_j\le11U/4\},\qquad L_U=\log U,
\]
where \(z_j=x_j+\ii y_j\) are the zeros of \(H_t\). The weight
\(\psi_U\) is supported on \([U/2,5U/2]\) and equals one on \([U,2U]\).
For an even convex \(g\in C^2([-1,1])\) with \(g(0)=0\), let
\[
 Q_{U,g}(t)=\sum_j\psi_U(x_j)g(y_j),\qquad
 \mathcal D_{U,g}(t)=\sum_{j:y_j\ne0}\psi_U(x_j)\frac{g'(y_j)}{y_j},
 \qquad G=\|g\|_\infty .
\]
Both quantities are nonnegative.

\begin{proposition}[Broad convex balance]
\label{gt-convex-balance}
The function \(Q_{U,g}\) is absolutely continuous and, with
\(D_\chi=C_{\rm u}^2(36\|\chi''\|_\infty+1200\|\chi'\|_\infty)\),
\begin{equation}
 Q_{U,g}(\beta)+\int_\alpha^\beta\mathcal D_{U,g}(t)\,\dd t
 \le Q_{U,g}(\alpha)+D_\chi G(\beta-\alpha)L_U^3
 \qquad(0\le\alpha\le\beta\le1/2).
 \label{gt-convex-integrated}
\end{equation}
\end{proposition}

\begin{proof}
For \(z=x+\ii y\) in the support strip, put \(w=x+4\ii\). Its
exterior force after removing the cluster \(\mathcal C_U\) is
\begin{equation}
 E_U(z)=\mathcal L_t(w)-\sum_{k\in\mathcal C_U}\frac1{w-z_k}
     +\sum_{k\notin\mathcal C_U}
        \left(\frac1{z-z_k}-\frac1{w-z_k}\right).
 \label{gt-broad-exterior}
\end{equation}
Indeed, start with the opposite-paired logarithmic derivative of
the even canonical product, as in the proof of
Proposition~\ref{gt-exterior-proposition}. At a zero use its regular
part after removing the self term. Subtract the value at \(w\) before
passing to the full exterior series. For \(k\notin\mathcal C_U\), the
horizontal distance is at least \(U/4\ge2\), so the difference in
\eqref{gt-broad-exterior} has modulus at most \(5|x-x_k|^{-2}\).
The complete unit bound gives at most
\(2C_{\rm u}[\log(2+x)+\log(h+2)]\) occurrences in a shell
\(h\le|x-x_k|<h+1\). Since \(\log(2+x)\le2L_U\), the difference
series over \(|x-x_k|\ge m\), for integers \(m\ge2\), has bound
\[
 \frac{10C_{\rm u}}{m-1}[2L_U+\log(m+1)+1].
\]
This follows by integrating the decreasing shell majorant, and
includes both horizontal tails. Taking \(m=\lfloor U/4\rfloor\)
bounds the full difference series by \(640C_{\rm u}L_U/U\),
using \(m-1\ge U/16\) and \(\log(m+1)+1\le2L_U\).

The cluster occupies at most \(3U\) closed unit intervals, each with
count at most \(2C_{\rm u}L_U\). For its finite sum at \(w\), each
distance shell contains at most \(4C_{\rm u}L_U\) occurrences
and each denominator has modulus at least \(\sqrt{h^2+9}\).
The maximum shell index is at most \(\lceil9U/4\rceil\le3U\).
Summing, and using \eqref{gt-safe-line} for the first term, gives
\[
 \begin{gathered}
 \sum_{k\in\mathcal C_U}|w-z_k|^{-1}
 \le4C_{\rm u}L_U[1/3+1+\log(3U)]\le12C_{\rm u}L_U^2,\\
 \#\mathcal C_U\le6C_{\rm u}UL_U,\qquad |E_U(z)|\le100C_{\rm u}L_U^2 .
 \end{gathered}
\]

At a simple zero the heat equation gives velocity \(H_t''/H_t'\),
twice the regular part of \(\mathcal L_t\). To differentiate
\(Q_{U,g}\), symmetrize the terms inside \(\mathcal C_U\).
For two distinct endpoints, let \(p=x_j-x_k\), \(q=y_j-y_k\),
and let \((x_s,y_s)\) be the chord from \(z_k\) to \(z_j\) with
\(0\le s\le1\). The force pairs the gradient difference of
\(f(x,y)=\psi_U(x)g(y)\) with \((p,-q)/(p^2+q^2)\). Integrating the
gradient difference along the chord cancels the mixed derivatives.
Thus
\[
 \begin{split}
 \mathcal K_{jk}
 &=\int_0^1
   \frac{\psi_U(x_s)g''(y_s)q^2-\psi_U''(x_s)g(y_s)p^2}
        {p^2+q^2}\,\dd s,\\
 Q_{U,g}'
 &=-2\sum_{j<k\in\mathcal C_U}\mathcal K_{jk}
   +2\sum_{j\,\text{in support}}
       [\psi_U'(x_j)g(y_j)\Re E_U(z_j)
        +\psi_U(x_j)g'(y_j)\Im E_U(z_j)] .
 \end{split}
\]
Only the contributing roots require velocities. Endpoints outside
the support have zero gradient, with zero kernel for coincident
inactive endpoints. The cluster is a partition of the force in this
identity at each time.

For each contributing conjugate pair \(x\pm\ii y\), \(y\ne0\), the chord
is vertical and contributes
\(-2\psi_U(x)[g'(y)-g'(-y)]/(2y)=-2\psi_U(x)g'(y)/y\)
to \(Q_{U,g}'\). This is precisely the negative of the two corresponding
terms in \(\mathcal D_{U,g}\). The other vertical interior terms
are nonpositive by convexity. For the exterior imaginary part,
the reciprocal terms converge absolutely by the square tail, and
the cluster removal preserves conjugation symmetry. For a conjugate
group with real part \(u\) and imaginary parts \(\pm v\), writing
\(s=x-u\), the negative imaginary part of the two reciprocals is
\[
 \frac{2y(s^2+y^2-v^2)}
      {[s^2+(y-v)^2][s^2+(y+v)^2]} .
\]
It has the sign of \(y\) by \(|s|\ge2,\ |v|\le1\). A real exterior
root has the same sign property, with expression \(y/(s^2+y^2)\).
Because \(yg'(y)\ge0\), the exterior vertical term is nonpositive.
The exterior horizontal term still uses the full convention in
\eqref{gt-broad-exterior}.

The horizontal cost inside the cluster is bounded by
\(\|\psi_U''\|_\infty G(\#\mathcal C_U)^2\), and the exterior
horizontal cost by
\(2\|\psi_U'\|_\infty G(6C_{\rm u}UL_U)(100C_{\rm u}L_U^2)\).
These bounds prove \(Q'_{U,g}+\mathcal D_{U,g}\le D_\chi GL_U^3\)
where the support roots are simple.

For assembly, fix \(U,g\). Local zero-free boundaries and Rouch\'e
continuity give continuity of the finite root multisets in a collar
around the support, and hence of \(Q_{U,g}\). The multiple-root locus
in the padded box \([U/2-1,5U/2+1]\times[-2,2]\) over the time slab
is compact. At a root of multiplicity \(m\) at \((t_0,z_0)\), the heat
equation gives the rescaled leading polynomial
\[
 |s|^{-m/2}H_{t_0+s}(z_0+\sqrt{|s|}\,\zeta)
 \ longrightarrow\ 
 \frac{H_{t_0}^{(m)}(z_0)}{m!}
     \exp(-\operatorname{sgn}(s)\partial_\zeta^2)\zeta^m
\]
on bounded \(\zeta\) sets, on each side separately. In the convergent
Taylor expansion, \(s^k(z-z_0)^j\) has weighted degree \(2k+j\),
and its coefficient is determined by the \((2k+j)\)-th spatial
derivative. Degrees below \(m\) vanish and degree \(m\) gives the
displayed polynomial. Its zeros are distinct by the Hermite polynomial
identity, with a rotation for \(s<0\)
\cite[Proposition~3.1(ii)]{polymath}. Rouch\'e on disjoint small circles
around these rescaled zeros, and on the original isolating disk,
shows that all \(m\) roots are simple for nearby nonzero \(s\).
Thus the multiple-root points in the compact locus are isolated and
there are only finitely many event times.

On the complementary intervals the absolute chord bound is
\(\max(\|\psi_U\|_\infty\|g''\|_\infty,\|\psi_U''\|_\infty G)\).
The active gradient and the field are bounded for fixed \(U,g\).
These give a finite two-sided derivative bound, so continuity across
the finite partition gives absolute continuity. With the summand set
to zero at \(y=0\), \(\mathcal D_{U,g}\) is Borel under local enumeration,
bounded by \(6C_{\rm u}UL_U\|g''\|_\infty\), and integrable. Integrating
the one-sided inequality proves \eqref{gt-convex-integrated}. Its
constant uses only \(G\), and the first derivative has no collision
atom.
\end{proof}

\subsection{Uniform density and transverse moments}

\begin{theorem}[Transverse density on the closed heat slab]
\label{gt-transverse-density}
For each fixed \(0<b<1/2\) there is \(B_b\) such that for \(U\ge8\),
\(0<v<1\), and \(0\le t\le1/2\),
\begin{equation}
 N_t^\perp(U,v):=\#\{j: U\le x_j\le2U,\ |y_j|>v\}
 \le B_b U L_U\e^{-bvL_U}.
 \label{gt-density-dyadic}
\end{equation}
In particular, for every integer \(k\ge1\),
\begin{equation}
 \sum_{U\le x_j\le2U}|y_j|^k
 \le \frac{k! B_b}{b^k}\frac{U}{L_U^{k-1}}.
 \label{gt-transverse-moments}
\end{equation}
\end{theorem}

\begin{proof}
At time zero, \(H_0(z)=\xi((1+\ii z)/2)/8\). A zero \(x+\ii y\)
corresponds to a zeta zero with real part \(1/2-y/2\) and
ordinate \(x/2\). Choose \(c_0\in(2b,1)\), and put
\(b_0=c_0/2\). The Selberg--Jutila bound as stated in Maples--Rodgers
\cite[Theorem~2.4 and \S4, (18), arXiv v2]{maplesrodgers}
bounds \(\#\{\beta_0>\sigma,\ 0<\gamma< V\}\) by
\(D_{c_0}V\log V\,\e^{-c_0(\sigma-1/2)\log V}\), with one constant
uniform in the threshold. On \([U/2,5U/2]\) the zeta ordinates
are strictly between \(0\) and \(2U\). Critical-line reflection
accounts for both signs of \(y\) with one doubling. Taking
\(V=2U\) and \(\sigma=1/2+v/2\) gives a bound
\(C_{c_0} U L_U\e^{-b_0vL_U}\) on the time-zero support population
with \(|y|>v\), where we can use \(C_{c_0}=8D_{c_0}\).

Use \(g_U(y)=\cosh(bL_Uy)-1\), which is convex and has sup norm
at most \(U^b\). Integrating its derivative against the initial
threshold count gives
\[
 Q_{U,g_U}(0)\le C_{c_0} U L_U
     \int_0^\infty b\sinh(bw)\e^{-b_0 w}\,\dd w
 =C_{c_0}\frac{b^2}{b_0^2-b^2}U L_U.
\]
The flux in \eqref{gt-convex-integrated} is at most
\(D_\chi U^bL_U^3/2\), and
\(U^{b-1}L_U^2\le4/[\e^2(1-b)^2]\). Consequently
\(Q_{U,g_U}(t)\le K_b U L_U\), where
\[
 K_b=C_{c_0}\frac{b^2}{b_0^2-b^2}
            +\frac{2D_\chi}{\e^2(1-b)^2}.
\]
For \(bvL_U\ge1\) each counted occurrence contributes at least
\(\kappa\e^{bvL_U}\) on the plateau, where
\(\kappa=(1-\e^{-1})^2/2\). For \(bvL_U<1\) use the
ordinary population bound \(6C_{\rm u}UL_U\). This proves
\eqref{gt-density-dyadic} with
\(B_b=\max\{K_b/\kappa,6\e C_{\rm u}\}\). The layer-cake formula
\(k\int_0^1 v^{k-1}N_t^\perp(U,v)\,\dd v\), bounded by the exponential
integral over \([0,\infty)\), proves \eqref{gt-transverse-moments}.
The strict transverse threshold includes \(|y|=1\) in this integral
through thresholds below one.
\end{proof}

The starting density is arithmetic, while the convex balance uses the
heat equation, strip and complete field. The latter remains available
for an occupation estimate with a different test
\cite{tnddensity,gttransverse}.

\subsection{The radial comparison and its center average}

Fix \(T\ge16\), \(L=\log T\), \(\ell=L^{-1}\), \(1\le a\le L\),
and \(3\le R\le T/2\). The scale \(\ell\) is frozen throughout
\([T,2T]\). At arbitrary real centers define the transverse sum
\begin{equation}
 \Xi_{a,R}(t,X)=\ell^2\!
    \sum_{\substack{\rho:\,|\rho-X|\ge a\ell\\
                             |\Re\rho-X|\le R}}
    \frac{(\Im\rho)^2}{|\rho-X|^4}.
 \label{gt-transverse-xi}
\end{equation}

\begin{proposition}[Radial comparison for an isolated pair]
\label{gt-radial-bridge}
Take a conjugation-invariant pair of occurrences \(z_\pm=c\pm\sqrt D\)
with \(c,D\in\R\), \(|D|\le d\ell^2\) and an exterior collar
\(|\rho-c|\ge a\ell\) for every exterior occurrence, where
\(0\le d<a^2\). Write \(w=\rho-c=x+\ii y\), \(r=|w|\) and
\(\delta=d/a^2\). The signed summand of \(M=-K\) is
\begin{equation}
 q(\rho)=\Re\frac1{w^2-D}
 =\frac1{r^2-D}
  -\frac{2y^2(r^2+D)}{(r^2-D)[(r^2-D)^2+4Dy^2]} .
 \label{gt-radial-identity}
\end{equation}
With \(S_R=\sum_{\rho\,\mathrm{exterior},\,|x|\le R}q(\rho)\),
\begin{equation}
 \ell^2 S_R\ge
 \ell^2\sum_{\substack{\rho\,\mathrm{exterior}\\|x|\le R}}
    \frac1{r^2-D}-C_\delta\Xi_{a,R}(t,c),\qquad
 C_\delta=\frac{2(1+\delta)}{(1-\delta)^3}.
 \label{gt-radial-loss}
\end{equation}
\end{proposition}

\begin{proof}
In \eqref{gt-exterior-k} the two factors multiply to \(w^2-D\).
The real numerator is \(r^2-D-2y^2\), and
\(|w^2-D|^2=(r^2-D)^2+4Dy^2\). Subtracting from \((r^2-D)^{-1}\)
gives the identity. For either sign of \(D\), the collar gives
\(r^2\pm D\ge(1-\delta)r^2>0\) and
\(|w^2-D|\ge(1-\delta)r^2\). Thus the correction is between zero
and \(C_\delta y^2/r^4\). At a double it is exactly \(2y^2/r^4\).
Both selected occurrences have radius \(\sqrt{|D|}<a\ell\).
Every nearby exterior occurrence enters \(\Xi_{a,R}(t,c)\),
including equality at either cutoff, and this proves the bound.
\end{proof}

\begin{theorem}[Center-averaged transverse correction]
\label{gt-averaged-transverse}
Uniformly in \(t,T,a,R\) in the stated domains, the mean of
\eqref{gt-transverse-xi} satisfies
\begin{equation}
 \frac1T\int_T^{2T}\Xi_{a,R}(t,X)\,\dd X\le\frac{C_s}{a^3},
 \qquad
 \frac1T\int_T^{2T}\Xi_{a,R}(t,X)\,\dd X\le\frac{C_0}{a}
       \quad\text{using only the complete count},
 \label{gt-xi-means}
\end{equation}
where \(C_s=11\pi B_b/(3b^2)\) and \(C_0=3\pi C_{\rm u}\).
\end{theorem}

\begin{proof}
For \(r_0>0\) define
\[
 I(y,r_0)=\int_{x^2+y^2\ge r_0^2}\frac{y^2}{(x^2+y^2)^2}\,\dd x .
\]
It is zero at \(y=0\), and for \(|y|\ge r_0\) integration over the
whole line gives \(\pi/(2|y|)\). For \(s=|y|/r_0\in(0,1)\),
integrating the two rays using
\(x/[2(x^2+y^2)]+\arctan(x/|y|)/(2|y|)\) gives
\[
 I(y,r_0)=\frac{\arcsin s-s\sqrt{1-s^2}}{r_0s}.
\]
The numerator's derivative is \(2s^2/\sqrt{1-s^2}\), so its ratio to
\(s^3\) equals \(2\int_0^1 v^2(1-s^2v^2)^{-1/2}\,\dd v\le\pi/2\),
the value at \(s=1\). The cases therefore give
\begin{equation}
 I(y,r_0)\le\frac\pi2\min\{r_0^{-1},y^2r_0^{-3}\}.
 \label{gt-one-zero-integral}
\end{equation}

Only zeros with real part in \(\mathcal B_T=[T/2,5T/2]\) enter the
center average. This band lies in the closed dyadic bands at scales
\(T/2,T,2T\). Since \(\log(T/2)\ge3L/4\), the sum of the three
moment allowances in \eqref{gt-transverse-moments} for \(k=2\) is
at most \((11/3)(2B_b/b^2)T/L\). The complete count over
\(\mathcal B_T\) is at most \(6C_{\rm u}TL\) by the unit bound.
With \(r_0=a\ell\), Tonelli and the substitution \(x=\Re\rho-X\)
give an upper bound \(\ell^2 T^{-1}\sum_{\Re\rho\in\mathcal B_T}
I(\Im\rho,r_0)\) by extending each center integral to the line
while retaining the radial cutoff. The two choices in
\eqref{gt-one-zero-integral} give the two stated constants.
\end{proof}

The powers in \eqref{gt-xi-means} refer to the same averaged observable.
For \(a=L^\eta\), \(0<\eta<1\), the two upper bounds are
\(O(L^{-3\eta})\) and \(O(L^{-\eta})\) while the physical radius
tends to zero. At fixed \(a\) both bounds are \(O(1)\).
For \(h>0\) the first upper bound gives Lebesgue measure
at most \(C_sT/(a^3h)\) for \(\Xi_{a,R}>h\) on \([T,2T]\).
The value at a designated center in \eqref{gt-radial-loss} still
requires a pairwise estimate. A collar of the chosen radius is an
independent premise of that application.

\subsection{Packing and exceptional interactions on a slice}

Now impose \(d\ge0\), \(\max(2,4\sqrt d)\le a\le L\), and
\(3\le R\le T/4\), and put \(r_0=a\ell\). At time \(t\), let
\(\mathcal P_t\) be the distinct unordered conjugation-invariant
pairs of zero occurrences in Proposition~\ref{gt-radial-bridge}
with \(c\in[T,2T]\) and the specified collar. A double counts as
one pair. The two occurrences of any qualifying pair lie within
\(\sqrt d\,\ell\le r_0/4\) of the center. If distinct pairs shared
one occurrence, their centers would lie within \(2\sqrt d\,\ell\)
of one another, and the other occurrence of the second pair would
be exterior to the first at distance at most \(3r_0/4\) from its
center. This contradicts the collar, and sharing both occurrences
gives the same unordered pair. A member of one pair is therefore
exterior to each other pair. It follows that
\begin{equation}
 |c_1-c_2|\ge(a-\sqrt d)\ell\ge3r_0/4,\qquad
 \#\mathcal P_t\le1+\frac{4TL}{3a}.
 \label{gt-pair-packing}
\end{equation}

\begin{theorem}[Negative interactions on a slice]
\label{gt-pair-negative-count}
Define \(N_-(t)=\#\{P\in\mathcal P_t:S_R(P)<0\}\) and
\(N_M(t)=\#\{P\in\mathcal P_t:M(P)<0\}\), where \(M(P)\) is the
complete exterior sum. For fixed \(0<b<1/2\),
\begin{equation}
 N_M(t)\le N_-(t)\le
 5B_b\left(\frac{31}{9}+\frac{8}{3ab}\right) TL\e^{-2ba/3}.
 \label{gt-negative-slice}
\end{equation}
\end{theorem}

\begin{proof}
The bound \(|w^2-D|\ge(15/16)|w|^2\) follows from
\(d/a^2\le1/16\). The unit count in each horizontal shell
\(n\le|x|<n+1\), \(n\ge1\), is
\(O(C_{\rm u}[\log(2+c)+\log(n+2)])\).
The resulting inverse-square tail is absolutely convergent,
and \(|x|<1\) contains finitely many exterior occurrences.
For the omitted terms \(|x|>R\) the numerator
\(x^2-y^2-D\) is positive, since \(|y|\le1\), \(|D|\le1/16\)
and \(R\ge3\). Thus \(M(P)\ge S_R(P)\).

If \(q(\rho)<0\), then \(y^2>(r^2-D)/2\ge15r_0^2/32\), and
hence \(|y|>v_0=2r_0/3\). Negativity also requires
\[
 |\Re\rho-c|<\sqrt{y^2+D}\le |y|+r_0/4 .
\]
By \eqref{gt-pair-packing}, such a window around a fixed occurrence
contains at most \(5/3+8|y|/(3r_0)\) eligible centers.
Each pair with \(S_R<0\) has at least one negative summand, and
every such occurrence has real part in \(\mathcal B_T\). Consequently
\begin{equation}
 N_-(t)\le
 \sum_{\substack{\Re\rho\in\mathcal B_T\\ |y|>v_0}}
        \left(\frac53+\frac{8|y|}{3r_0}\right).
 \label{gt-negative-incidences}
\end{equation}
This counts an occurrence at each center it can affect. Multiple
incidences for one pair only increase the upper bound.

For the three dyadic bands covering \(\mathcal B_T\), the ratios of
\(U\log U\,\e^{-bv\log U}\) to \(TL\e^{-bvL}\) are bounded by
\(2^{b-1},1,3\), respectively. Thus
\[
 N_t^\perp(v;\mathcal B_T):=\#\{\rho:\Re\rho\in\mathcal B_T,\ |y|>v\}
 \le5B_b TL\e^{-bvL}\qquad(0<v<1).
\]
Layer cake gives the first moment over the strict threshold:
\[
 \begin{split}
 \sum_{\substack{\Re\rho\in\mathcal B_T\\ |y|>v_0}}|y|
 &=v_0N_t^\perp(v_0;\mathcal B_T)+\int_{v_0}^1N_t^\perp(v;\mathcal B_T)\,\dd v\\
 &\le5B_b TL\e^{-bv_0L}\left(v_0+\frac1{bL}\right).
 \end{split}
\]
Substituting these two bounds into \eqref{gt-negative-incidences}
gives \eqref{gt-negative-slice}.
\end{proof}

\paragraph{A large transverse loss at a qualifying center.}
The intervals \(I_c=[c-r_0/4,c+r_0/4]\) at the qualifying centers
are disjoint by \eqref{gt-pair-packing}. Moving from \(c\) to
\(X\in I_c\) shifts distances by at most \(r_0/4\).
Every occurrence counted in \(\Xi_{a,R}(t,c)\) then has radial
distance at least \(3r_0/4\) and horizontal distance at most
\(R'=R+r_0/4\). Its new distance is no more than \(5/4\) times its
old one. With \(\theta_*=(4/5)^4\), this gives
\[
 \Xi_{a/2,R'}(t,X)\ge\theta_*\Xi_{a,R}(t,c)\qquad(X\in I_c).
\]
All these intervals lie in \(J=[T-r_0/4,2T+r_0/4]\).
For \(X\in J\) the enlarged horizontal cutoff places the contributing
zeros in \([T-R-r_0/2,2T+R+r_0/2]\subseteq\mathcal B_T\), using
\(R\le T/4\), \(r_0\le1\) and \(T\ge16\).
Apply the same one-zero integral over \(J\) with radius \(r_0/2\).
By \eqref{gt-one-zero-integral} and the population bounds used in
Theorem~\ref{gt-averaged-transverse}, the integral is at most
\(8C_sT/a^3\), or \(2C_0T/a\) using only the count.
Each qualifying center with \(\Xi_{a,R}(t,c)\ge h>0\) contributes
at least \(\theta_*h r_0/2\) on its interval. Thus the number
of such centers is bounded by
\begin{equation}
 \frac{16C_s}{\theta_*}\frac{TL}{a^4h},\qquad
 \frac{4C_0}{\theta_*}\frac{TL}{a^2h}
     \quad\text{using only the complete count}.
 \label{gt-xi-exceptions}
\end{equation}
Both allowances can be intersected with \eqref{gt-pair-packing}.
Since \(C_\delta\le8704/3375\) in this class, either
\(\ell^2 S_R\le-\epsilon\) or \(\ell^2M\le-\epsilon\) implies
\(\Xi_{a,R}(t,c)\ge(3375/8704)\epsilon\), by
\eqref{gt-radial-loss} and positivity of the tail.

There is also a pointwise count-only bound
\(\Xi_{a,R}(t,c)=O(L/a^2+1/L)\). For \(|x|\le1\) two unit windows
supply at most \(4C_{\rm u}L\) terms, each at most \(a^{-2}\).
In the other shells each term is at most \(\ell^2 n^{-4}\) by the
strip, and summing the shell counts from the proof of
Theorem~\ref{gt-pair-negative-count} gives \(O(1/L)\).
Thus for a fixed threshold \(h\) and \(a/\sqrt L\to\infty\), the
pointwise bound already excludes \(\Xi_{a,R}\ge h\) at sufficiently
large height. A growing-radius example of the two-power gain in
\eqref{gt-xi-exceptions} is \(a=L^\eta\) with \(0<\eta<1/2\).
The sign count \eqref{gt-negative-slice} includes arbitrarily small
negative interactions. Its allowance divided by \(TL\) decays
exponentially in \(a\), while the count-only packing allowance
has order \(1/a\) for \(a=L^\eta\), \(0<\eta<1\). At \(a=L\)
the signed upper allowance still grows like \(T^{1-2b/3}L\).
These are bounds with a denominator set by the height band.
A fraction among qualifying pairs would require a lower bound
on their number.

\subsection{Pair-time occupation at a fixed collar}

\begin{theorem}[Time occupation of negative interactions]
\label{gt-pair-occupation}
Assume \(T\ge64\) and the same slice class as in
Theorem~\ref{gt-pair-negative-count}, with \(a,d,R,r_0\)
fixed in time. For \(0\le\alpha\le\beta\le1/2\) and
\(\tau=\beta-\alpha\),
\begin{equation}
 \begin{split}
 \int_\alpha^\beta N_M(t)\,\dd t
 &\le\int_\alpha^\beta N_-(t)\,\dd t\le
 B_b\left(\frac{93}{b^2}+\frac{279}{ab^3}\right)\frac{T}{L}\
 +\frac{32}{3}D_\chi
    \left(\frac12+\frac{L}{3a}\right)\tau L^3 .
 \end{split}
 \label{gt-occupation-bound}
\end{equation}
The upper bound can also be intersected with \(\tau\) times each
slice allowance in \eqref{gt-pair-packing} and
\eqref{gt-negative-slice}.
\end{theorem}

\begin{proof}
Use
\[
 g(y)=\frac{y^2}{2}+\frac{|y|^3}{3r_0},
 \qquad \frac{g'(y)}{y}=1+\frac{|y|}{r_0}\ (y\ne0),\qquad
 G=\frac12+\frac{L}{3a}.
\]
The function is even, \(C^2\) and convex. Let
\(A_T=\psi_{T/2}+\psi_T+\psi_{2T}\) and write \(Q_A,\mathcal D_A\)
for the corresponding sums of the three observables. The three
plateaus cover \(\mathcal B_T\), so \(A_T\ge1\) there, and \(A_T\le3\)
with support in \([T/4,5T]\). By \eqref{gt-negative-incidences}
and the choice of \(g\), \(N_-(t)\le(8/3)\mathcal D_A(t)\).

The five dyadic bands at scales \(T/4,T/2,T,2T,4T\) cover
\([T/4,5T]\). Their scales sum to \(31T/4\) and each logarithm is
at least \(2L/3\). For \(k=2,3\), equation
\eqref{gt-transverse-moments} then bounds the sums of \(|y|^k\)
on the support by \(93B_bT/(4b^2 L)\) and
\(837B_bT/(8b^3L^2)\), respectively. Inserting the coefficients of
\(g\) and the bound \(A_T\le3\) gives
\[
 Q_A(\alpha)\le\frac{B_b}{8}
     \left(\frac{279}{b^2}+\frac{837}{ab^3}\right)\frac{T}{L}.
\]
Summing \eqref{gt-convex-integrated} over the three profiles and
dropping \(Q_A(\beta)\ge0\) bounds the integral of
\(\mathcal D_A\) by this allowance plus \(4D_\chi G\tau L^3\),
since the three log cubes sum to \(3L^3+6L(\log2)^2\le4L^3\).
Multiplying by \(8/3\) gives \eqref{gt-occupation-bound}.

To justify the integrations, the slice conditions for the targets,
collar, and \(S_R\) use a fixed padded finite band of roots.
Off its finite multiple-event time set, local simple-root charts
make the pair enumerations and all inequalities Borel. At event
times the counts are defined directly on the finite multiset.
Exhausting the full divisor by larger bounded bands gives \(M\)
as a measurable limit on qualifying pairs, with a uniform
inverse-square tail for \(c\in[T,2T]\). The charts are used only
for measurability. Membership in \(\mathcal P_t\) is tested at
each slice, including at \(D=0\).
\end{proof}

For fixed \(a,b\) and fixed positive \(\tau\), the occupation
allowance divided by \(TL\) is
\(O_{a,b}(L^{-2})+O_{a,\chi}(L^3/T)\).
For variable positive \(\tau\), division by \(\tau TL\) gives
the minimum of an \(O_{a,b}(1)\) slice allowance and the
common-time allowance
\(O_{a,b}((\tau L^2)^{-1})+O_{a,\chi}(L^3/T)\).
The duration condition for decay of the latter bound at fixed
\(a,b\) is \(\tau L^2\to\infty\). On a mean-spacing squared
time scale the starting-moment allowance can still pay the entire
slice scale. An exclusion for a particular pair would require a
lower bound for its time in the counted class with negative
interaction. At large height the allowance leaves room even for
one pair on the entire slab. A signed transport bound also uses an
independently supplied incoming value.

\paragraph{Source dependence and matched comparison.}
The same estimates are available to direct evaluation given the
complete divisor, strip, density and pair geometry. The occupation
gain uses the integrated heat identity as well. For a fixed-duration
comparison, integrating the slice count gives only the
\(O_{a,b}(1)\) normalized allowance. No evaluator-specific
information enters that comparison.

The opposite-sign fixed-slice examples from Section~\ref{sec:gt-height}
can also obey the upper transverse density and moment bounds.
The real insertions add zero transverse moment. For the nonreal
insertion at \(T=c\), the two extra nonreal occurrences at positive
height have \(|y|=2\sqrt3/\log c\). On every positive dyadic band
\([U,2U]\), \(U\ge8\), entered by this pair, \(\log U\le\log(2c)\).
Whenever the strict density threshold counts the pair,
\(v\log U\) is therefore uniformly bounded. The two occurrences
are absorbed by an enlarged density constant, and deleting
occurrences reduces all nonnegative counts and moments. Their
near sums \(S_3\) have exact opposite signs at the designated
doubles with \(a=2\); the full interactions \(M\) have opposite
signs for all sufficiently large \(c\). Along the unbounded choices
with \(H_0(c)\ne0\), the modified functions differ from the original
theta field. No common heat-time family is supplied by this
fixed-slice construction.
Quadratic and quartic heat controls can share the conjugate
dissipation but lack the canonical growing count law. The
weighted-square and finite-source controls change the full
reflected theta expansion, and fractional profiles lack the
integer divisor required by the incidence count. The separate
premise screens specify the scope of the remaining controls
\cite{gtcounts,gttransverse}.
